\par\Needspace{12\baselineskip}\bigskip
\pdfbookmark[2]{Exposition, Second Section}{pallas-II-section2}
\begin{center}
{\large\bfseries Second Section.\par}
\medskip
\textit{Instantaneous variations of the elements produced by perturbations.}
\end{center}
\pdfbookmark[3]{Exposition, 7}{pallas-II-article7}
\begin{center}7.\end{center}
\noindent\hspace*{\parindent}The elliptic motion that we have just explained presupposes that the mutual attractions of the Sun and the planet are the only forces acting. If still other forces have to be taken into account, the motion of the planet about the Sun becomes more complicated. But whatever that motion may be, it is clear that if these forces were suddenly extinguished, from that moment the planet would move along a conic section whose elements would be functions of the quantities that determine the position and the velocity acquired by the planet at that instant, that is, of the six quantities \(x,y,z,\dfrac{\mathrm{d}x}{\mathrm{d}t}\mathpunct{\raisebox{.45ex}{,}}\dfrac{\mathrm{d}y}{\mathrm{d}t}\mathpunct{\raisebox{.45ex}{,}}\dfrac{\mathrm{d}z}{\mathrm{d}t}\mathpunct{\cdot}\) We confine ourselves to the case in which this conic section is always an ellipse, \SourceLeaf{456}and we shall call it the \textls[80]{osculating ellipse}. We thus see that all the finite equations between the elements and the quantities \(x,y,z,r,\dfrac{\mathrm{d}x}{\mathrm{d}t}\mathpunct{\raisebox{.45ex}{,}}\dfrac{\mathrm{d}y}{\mathrm{d}t}\mathpunct{\raisebox{.45ex}{,}}\dfrac{\mathrm{d}z}{\mathrm{d}t}\mathpunct{\raisebox{.45ex}{,}}\dfrac{\mathrm{d}r}{\mathrm{d}t}\mathpunct{\raisebox{.45ex}{,}}\) that we developed for purely elliptic motion, or more precisely all those that hold in that case, will also hold for the osculating ellipse. But since the perturbing forces continue to act without interruption, the osculating ellipse will vary from one instant to the next. Its elements will be variable quantities. If, however, we succeed in assigning numerical values to all the elements of the osculating ellipse at each instant, the problem of determining the motion will be solved. We shall use the same formulas as in the case of unperturbed motion to calculate the position of the planet as well as its velocity and the direction of its motion. This investigation reduces to two problems. First we shall determine the instantaneous variations undergone by the several elements under the action of the perturbing forces. Then we must integrate these differential expressions to obtain the values of the variable elements at each instant.
